\documentclass[11pt]{article}
\usepackage{amsmath,amssymb,amsthm,hyperref}
\newtheorem{lemma}{Lemma}
\newcommand{\E}{\mathbb E}
\title{Forced-success mixture: concentration and scaled Laguerre tests}
\author{Independent companion calculation for the power lower-bound argument}
\date{30 September 2026}
\begin{document}\maketitle
Let $p_1,\ldots,p_m$ be a monotone uniform-interval comparison tensor
with row weights $w$. Write $\delta_j=1-p_j$, $s=\sum_j\delta_j$,
$\ell=\lfloor m/4\rfloor$, $N=m-\ell$, and $\alpha=\ell/m$.
Take $m$ sufficiently large, so $1/5\le\alpha\le1/4$.
Let $B$ be a uniformly chosen $\ell$-subset of the columns, $R=B^c$,
and define
\[
 A_B=\prod_{j\in B}p_j,\qquad
 \overline A=\E_B A_B,\qquad
 x=\alpha s,\qquad x_R=\frac\ell N\sum_{j\in R}\delta_j.
\]
The finite mean-deficit balance theorem in
\path{../rational_designs/asymptotic/finite_endpoint_odds_balance.tex}
implies the global estimate
\[
 \max_j\delta_j\le\frac{256(s+1)}m.                 \tag{1}
\]
For $s\le m/256$ this follows from that theorem's odds bound; for
larger $s$ it follows from $\delta_j\le1$.

\begin{lemma}[Laplace mass]
For every fixed $c>0$,
$\E_w e^{-cs}\le C_c/m$.
Consequently $\E_w(1+s)^a e^{-cs}\le C_{a,c}/m$ for every fixed
nonnegative integer $a$.
\end{lemma}
\begin{proof}
It suffices to treat $0<c\le1$. On $s\le m/1024$, (1) gives
$\max_j\delta_j\le1/2$ for sufficiently large $m$. Choose any
$k=\lfloor cm/1024\rfloor$ columns. Their success product is at least
\[
 \exp\left(-2\sum_{j\in I}\delta_j\right)
 \ge\exp\left(-512k(s+1)/m\right)
 \ge e^{-c(s+1)/2}.
\]
Thus $e^{-cs}$ on this region is at most $e^{c/2}$ times that
success product. Its expectation is exactly $1/(k+1)$.
On the complementary region $e^{-cs}\le e^{-cm/1024}$.
These two bounds prove the first assertion, after enlarging the
constant to cover bounded $m$. The second follows by absorbing the
polynomial into $e^{-cs/2}$.
\end{proof}

\begin{lemma}[Weighted mixture concentration]\label{mix}
At every row,
\[
 \E_B x_R=x,\qquad
 \E_B(x_R-x)^2\le C\frac{(s+1)^2}{m},\qquad
 \E_B[A_B|x_R-x|]\le C e^{-\alpha s/2}\frac{s+1}{\sqrt m}.
\tag{2}
\]
If a differentiable function $P$ satisfies
\[
 |P'(u)|\le H(1+u)^a e^{\theta u}\quad(u\ge0),
 \qquad 0<\theta\le1/16,
\]
then
\[
 (\ell+1)\left|\E_w\E_B[A_BP(x_R)]-
                  \E_w[\overline A P(x)]\right|
 \le C_a Hm^{-1/2}.                               \tag{3}
\]
\end{lemma}
\begin{proof}
Sampling without replacement gives the exact variance
\[
 \operatorname{Var}_B(x_R)=
 \frac{\ell^3}{N(m-1)}
 \left(\frac1m\sum_j\delta_j^2-\frac{s^2}{m^2}\right).
\]
Since $\sum_j\delta_j^2\le256s(s+1)/m$, this proves the first two
assertions in (2). Maclaurin's inequality gives
$\overline A\le(1-s/m)^\ell\le e^{-\alpha s}$.
As $A_B^2\le A_B$, Cauchy--Schwarz gives the last assertion in (2).
Also $\max\{x,x_R\}\le\alpha s/(1-\alpha)\le s/3$.
The mean value theorem, (2), and $\alpha\ge1/5$ therefore bound the
rowwise difference in (3) by
\[
 C_a Hm^{-1/2}(1+s)^{a+1}
 e^{-\alpha s/2+\theta s/3}
 \le C_a Hm^{-1/2}(1+s)^{a+1}e^{-s/16}.
\]
The Laplace-mass lemma and $\ell+1=O(m)$ prove (3).
\end{proof}

Write $L_j(y)=\sum_{k=0}^j(-1)^k\binom jk y^k/k!$ for the
Laguerre polynomials, and put
\[
 q_d(y)=\sum_{j=0}^{d-1}(1-j/d)L_j(y),\qquad
 R_d(x)=\left(\frac2{\theta d}-x\right)q_d(\theta x)^2.
\]
The standard Laguerre estimates $|L_j^{(k)}(y)|\le\binom jk e^{y/2}$
for $y\ge0$ imply
\[
 |R_d'(x)|\le C_\theta d^3(1+x)e^{\theta x}.
\tag{4}
\]
For $h=\lfloor d/8\rfloor$ and
$P_d(x)=h^{-2}(\sum_{j<h}L_j(\theta x))^2$ they similarly give
$|P_d'(x)|\le C_\theta d e^{\theta x}$.
Thus (3) yields mixture-comparison errors $O(d^3/\sqrt m)$ and
$O(d/\sqrt m)$ respectively.

\begin{lemma}[Scaled taper under the exact beta law]
Let $\beta_\ell$ have density
\[
 b_\ell(x)=\frac{\ell+1}{\ell}(1-x/\ell)^\ell
             \boldsymbol1_{[0,\ell]}(x).
\]
For fixed $0<\theta\le1/16$, sufficiently large $d$ and $\ell$,
\[
 \int R_d\,d\beta_\ell\ge c_\theta>0,
 \qquad
 \int P_d\,d\beta_\ell\le C_\theta/d.
\tag{5}
\]
Moreover $P_d(x)\ge1/4$ whenever $0\le x\le2/(\theta d)$.
\end{lemma}
\begin{proof}
The ratio $b_\ell(x)/e^{-\theta x}$, extended by zero beyond $\ell$,
is decreasing, since its logarithmic derivative on the interior is
$-1/(1-x/\ell)+\theta<0$. Reweighting the probability measure
proportional to $q_d(\theta x)^2e^{-\theta x}dx$ by this decreasing
ratio decreases its mean. The exact Laguerre identities give
\[
 \frac{\int_0^\infty xq_d(\theta x)^2e^{-\theta x}dx}
      {\int_0^\infty q_d(\theta x)^2e^{-\theta x}dx}
 =\frac3{\theta(2d+1)}.
\]
It follows that
\[
 \int R_d\,d\beta_\ell
 \ge\frac1{2\theta d}\int q_d(\theta x)^2\,d\beta_\ell.
\]
On $0\le x\le1/(8\theta d)$, each $L_j(\theta x)$ with $j<d$
is at least $1/2$: its coefficient formula bounds the deviation from
one by $e^{j\theta x}-1\le e^{1/8}-1<1/2$.
Therefore $q_d(\theta x)\ge d/4$. For large $d,\ell$, the beta
density is bounded below by an absolute constant on this interval.
The last integral is at least $c_\theta d$, proving the first claim.

For the second claim, $b_\ell(x)\le2e^{-\theta x}$ and Laguerre
orthogonality give
$\int P_d\,d\beta_\ell\le2/(\theta h)=O_\theta(d^{-1})$.
If $x\le2/(\theta d)$ and $j<h\le d/8$, the same coefficient estimate
gives $|L_j(\theta x)-1|\le e^{1/4}-1<1/2$. Thus the sum defining
$P_d$ is at least $h/2$, as claimed.
\end{proof}

The exact beta law arises from the identity
\[
 (\ell+1)\E_w[A_B\prod_{j\in I}\ell(1-p_j)]
 =\int x^{|I|}\,d\beta_\ell(x),\qquad I\subseteq B^c.
\]
Consequently expectations of symmetric multiaffine polarizations on
$R$ are exact beta integrals. The comparison between a polarization
and its scalar value at $x_R$ is a separate required input; it is
not proved in this companion note. Thus this note alone does not
claim a power lower bound for the number of rows.
\end{document}
